\documentclass[10pt,letterpaper]{article}
\usepackage[T1]{fontenc}
\usepackage{lmodern,amsmath,amssymb,booktabs,array,geometry,xcolor}
\usepackage[hidelinks]{hyperref}
\geometry{margin=0.76in,headheight=14pt,headsep=15pt}
\definecolor{ink}{HTML}{173A51}
\makeatletter
\def\ps@draft{\def\@oddhead{\small\color{ink}Spin(10) and global PQ: discussion draft\hfil 14 October 2026}
\def\@evenhead{\@oddhead}
\def\@oddfoot{\small Discussion draft; bounded results and open physical questions\hfil\thepage}
\def\@evenfoot{\@oddfoot}}
\makeatother
\pagestyle{draft}
\setlength{\parindent}{0pt}\setlength{\parskip}{6pt}
\newcommand{\PQ}{\mathrm{PQ}}
\newcommand{\Spin}{\operatorname{Spin}}
\newcommand{\Tr}{\operatorname{Tr}}
\newcommand{\hc}{\mathrm{h.c.}}
\newcommand{\pathref}[1]{\texttt{\detokenize{#1}}}
\newcommand{\sect}[1]{\vspace{5pt}{\large\bfseries\color{ink}#1}\par}
\begin{document}
\thispagestyle{draft}
{\LARGE\bfseries\color{ink} An AI-assisted investigation of a\par
Spin(10) $\times$ global U(1)$_{\PQ}$\par
gauge/scalar candidate\par}
\vspace{4pt}
{\large Construction, bounded results, and failed tests}\par
{\small Solo research discussion note, prepared with AI assistance. Version 2.}

\sect{Abstract}
This note defines a nonsupersymmetric, renormalizable gauge/scalar construction
and summarizes selected calculations already recorded for it. A complete
scalar-invariant basis and a locally positive high-scale scalar point have been
obtained. An initial one-light-Higgs tuning fails its heavy-relaxed tree-level
quartic test; a later distinct point has a positive relaxed quartic. The
recorded gauge-matching equations also retain a nonzero mismatch; model-specific
finite two-loop heavy thresholds remain incomplete. Selected classical symmetry
identities and one-loop scalar operator closure have been checked within stated
scopes. These results neither establish a physically viable grand unified theory
nor exclude the entire construction. This is a technical discussion draft, not
a new scientific acceptance decision.

\sect{1. Assumptions and fields}
Work in four spacetime dimensions with gauge group $\Spin(10)$ and global PQ
selection rules. The simply connected group permits the fermionic spinor $16$;
local tensor calculations use the $\mathfrak{so}(10)$ action. A possible quotient
by common gauge/PQ centers has not been fixed. No supersymmetry, explicit soft
PQ violation, or additional scalar representation is assumed.

\begin{center}
\begin{tabular}{lccc}
\toprule Field & Representation/type & $q_{\PQ}$ & Real scalar coordinates\\
\midrule
$\Phi_{ij}$ & real symmetric traceless $54$ & $0$ & $54$\\
$\Sigma_{ijklm}$ & complex self-dual $126$ & $+2$ & $252$\\
$\phi_i$ & complex $10$ & $-2$ & $20$\\
$S$ & complex singlet & $-4$ & $2$\\
$\psi_p$, $p=1,2,3$ & left-handed $16_F$ & $+1$ & ---\\
\bottomrule
\end{tabular}
\end{center}
Indices run from $0$ to $9$, with Euclidean vector metric $\delta_{ij}$ and
$\epsilon_{0\cdots9}=+1$. Fix $*\Sigma=+i\Sigma$; its conjugate is the
$\overline{126}$. The scalar kinetic convention is
\[
\mathcal L_{\rm kin}=
\frac12(D_\mu\Phi_{ij})^*(D^\mu\Phi_{ij})
+\frac{1}{2\,5!}(D_\mu\Sigma_{ijklm})^*(D^\mu\Sigma_{ijklm})
+(D_\mu\phi_i)^*(D^\mu\phi_i)+(\partial_\mu S)^*(\partial^\mu S).
\]
The self-dual factor avoids double counting complementary five-index sectors.
The physical construction uses the usual gauge-curvature kinetic term and
$\mathcal L=\mathcal L_{\rm gauge}+\mathcal L_{\rm kin}+\mathcal L_Y-V$.
The allowed symbolic Yukawa structures are $\psi\psi\phi$ and
$\psi\psi\Sigma^*$, with complex symmetric family matrices and isometrically
normalized intertwiners. Their numerical matrices and right-handed-neutrino
spectrum have not been selected.

The field-content motivation follows Babu and Khan [1]; PQ/axion context is
discussed in [2,3]. The potential below is a separately specified source-derived
variant. Published scalar benchmarks, fermion fits, and unification conclusions
are not imported as results for it. No novelty claim is made.

\newpage
\sect{2. Complete frozen scalar definition}
Write $V=V_2+V_3+V_4$. Repeated tensor indices are summed over all ten values.
Define
\[
\begin{gathered}
N=\frac1{5!}\Sigma_{abcde}\Sigma^*_{abcde},\quad
u=\phi_i^*\phi_i,\quad t=\phi_i\phi_i,\quad
p_n=\Tr(\Phi^n),\ n=2,3,4,\\
K_{ij}=\frac1{4!}\Sigma_{iabcd}\Sigma^*_{jabcd},\qquad
T_{ij}=\frac1{4!}\Sigma_{iabcd}\Sigma_{jabcd},\\
L_{\Phi\Sigma}=\frac1{2\,3!}\Phi_{ij}\Phi_{kl}
\Sigma_{ikabc}\Sigma^*_{jlabc},\qquad
E=\frac1{2!\,2!\,3!}\Sigma_{ijlmn}\Sigma_{ijpqr}\Sigma^*_{lmpqr}\phi_n .
\end{gathered}
\]
$T$ is symmetric traceless. For increasing tuples $I,J$ of length $5-k$,
let $P_k[I,J]=\sum_{K\in\binom{\{0,\ldots,9\}}k}\Sigma_{KI}\Sigma_{KJ}$
and $Q_k=\sum_{I,J}|P_k[I,J]|^2$, for $k=0,1,2$.
Antisymmetric components with repeated indices vanish. If $A,B$ are disjoint,
$\varepsilon(A,B)$ is the parity of their concatenation relative to increasing
order; otherwise it is zero. The fourth real five-form quartic is
\[
\begin{split}
X_{131}={}&\sum_{I,J\in\binom{\{0,\ldots,9\}}4}
\ \sum_{\substack{A\subset I,\ |A|=3\\B\subset J,\ |B|=1}}
\varepsilon(A,I\setminus A)\varepsilon(B,J\setminus B)
\varepsilon(A,B)\varepsilon(I\setminus A,J\setminus B)\\[-2pt]
&\hspace{25pt}\times P_1[I,J]\,
\overline{P_1[\operatorname{sort}(A\cup B),
\operatorname{sort}((I\setminus A)\cup(J\setminus B))]} .
\end{split}
\]
The potential in these normalizations is
\begin{align*}
V_2={}&m_\Phi^2p_2+m_\Sigma^2N+m_\phi^2u+m_S^2|S|^2,\\
V_3={}&\mu_\Phi p_3+\mu_{\Phi\phi}\Phi_{ij}\phi_i^*\phi_j
+[z_6tS^*+\hc],\\
V_4={}&\lambda_{\Phi1}p_2^2+\lambda_{\Phi2}p_4
+\lambda_{\Phi\Sigma1}p_2N+\lambda_{\Phi\Sigma2}L_{\Phi\Sigma}\\
&+\lambda_{\Phi\phi1}p_2u+\lambda_{\Phi\phi2}\phi_i^*\Phi_{ij}\Phi_{jk}\phi_k
+\lambda_{\Phi S}p_2|S|^2\\
&+\lambda_{\Sigma1}Q_0+\lambda_{\Sigma2}Q_1
+\lambda_{\Sigma3}Q_2+\lambda_{\Sigma4}X_{131}\\
&+\lambda_{\Sigma\phi1}Nu+\lambda_{\Sigma\phi2}\phi_i^*K_{ij}\phi_j
+\lambda_{\phi1}u^2+\lambda_{\phi2}|t|^2\\
&+\lambda_{\Sigma S}N|S|^2+\lambda_{\phi S}u|S|^2+\lambda_S|S|^4\\
&+[z_4\Phi_{ij}T_{ij}S+z_K\Phi_{ij}\phi_i\phi_jS^*
+z_\eta E+z_D T_{ij}\phi_i\phi_j+\hc].
\end{align*}
All $m^2,\mu,\lambda$ coefficients are real; $z_6,z_4,z_K,z_\eta,z_D$
are complex. The masses have dimension two, the cubics dimension one, and
the quartics dimension zero. Hermitian conjugation includes the coefficient.
The shorthand $\mu_{\Phi\phi},\lambda_{\phi1,2},\lambda_{\phi S},z_\eta$
corresponds to repository names \pathref{muPhiPhi}, \pathref{lambdaPhiVector1,2},
\pathref{lambdaVectorS}, and \pathref{zEta}.

There are 29 Hermitian coefficient groups: four quadratic, three cubic, and
22 quartic. Their 24 real and five complex coefficients give 34 real directions
before rephasings. Independent $D_5$ character multiplicities give 34 complex
singlet slots before Hermitian pairing; explicit contractions attain that upper
bound. The generic scalar phase quotient leaves three relative coupling phases.
This basis result fixes the action, not a viable vacuum.

\newpage
\sect{3. Scalar results: a failed tuning and a later positive point}
The candidate high-scale background is
\[
\langle\Phi\rangle=\frac{\omega}{\sqrt{60}}
\operatorname{diag}(-2\mathbf1_6,3\mathbf1_4),\quad
\langle\Sigma\rangle=\frac{\sigma}{4\sqrt2}
\bigwedge_{k=0}^{4}(e_{2k+1}+i e_{2k}),\quad
\langle S\rangle=\frac v{\sqrt2},\quad\langle\phi\rangle=0.
\]
The first direction has the intended Pati--Salam/discrete-parity stabilizer;
the combined gauge orbit has 33 broken and 12 unbroken directions, consistent
with the SM Lie algebra. Global topology and the PQ quotient are separate
questions. Tadpoles and Hessians are derived from the same frozen potential.

At the recorded Stage 1 point, $\omega=\sqrt{60}$ in calculation units and
$\sigma/\omega=v/\omega=0.1$. The exact-decimal free coefficients and
tadpole-eliminated masses are specified in \pathref{STAGE1_POINT.json} and its
replay program. All 35 SM-irrep multiplicity blocks were evaluated. The real
Hessian has rank 294 and nullity 34: 33 gauge directions plus one PQ direction.
The remaining eigenvalues are positive in the canonical kinetic metric;
the smallest recorded physical mass squared is approximately $0.08656$.
This is local positivity modulo symmetry, not global boundedness or a
phenomenologically fitted scale. This point has no light Higgs doublet.

The distinct Stage 2 point changes
$\operatorname{Re}z_6,\operatorname{Re}z_K,\operatorname{Re}z_\eta$ to $0.3$
in the stored dimensionless convention and tunes
$m_\phi^2/\omega^2\simeq-0.1337864654$ to the exact algebraic root recorded in
\pathref{STAGE2_POINT.json}. It has rank 290 and nullity 38, including four
real coordinates of one massless complex SM Higgs doublet. All physical heavy
quadratic modes remain positive. This statement alone is insufficient for
local stability along the tuned flat direction.

\sect{Decisive additional test}
For a canonically normalized neutral field $h$, $H^\dagger H=h^2/2$, expand
around the tuned point with physical heavy coordinates $X_A$:
\[
V=V_0+a_{\rm direct}h^4+h^2j_AX_A
+\tfrac12X_A\mathcal H_{AB}X_B+\cdots .
\]
Minimization on the positive heavy subspace gives
$X=-\mathcal H^+j h^2+\cdots$, and hence
\[
\lambda_{\rm eff}=4\left(a_{\rm direct}
-\tfrac12j^T\mathcal H^+j\right),\qquad
V_{\rm eff}=\lambda_{\rm eff}(H^\dagger H)^2+\cdots .
\]
The recorded parent-action evaluation gives
\begin{center}
\begin{tabular}{lr}
\toprule Quantity in $(H^\dagger H)^2$ normalization & Value\\
\midrule Fixed-heavy direct quartic & $+0.19910$\\
Heavy relaxation subtraction & $56.68893$\\
Relaxed effective quartic & $\boldsymbol{-56.48984}$\\
\bottomrule
\end{tabular}
\end{center}
The negative sign makes the Stage 2 point a failed local scalar minimum at
tree level. It does not invalidate the separate Stage 1 result or prove that
every parameter point fails. A direct parent-potential expansion along the
heavy-field valley reproduces the quartic; heavy-coordinate shear and complex
Higgs-phase controls preserve it. These extraction routes share the frozen
tensor implementation. No global-minimum or loop-corrected stabilization claim
is made.

\sect{A distinct positive-quartic point}
The later recorded point in \pathref{canonical_so10_positive_higgs/POINT.json}
retains the Stage 1 VEVs but uses $z_6/\omega=z_K=0.03$, $z_\eta=0.3$,
$\lambda_{\Phi\phi1}=-0.002$ and an exact algebraic tuning
$m_\phi^2/\omega^2\simeq0.007302047861$. Its real Hessian has rank 290 and
nullity 38, with positive heavy modes. The direct quartic is $0.19845235$;
heavy relaxation subtracts $0.09359411$, leaving
\[
\boxed{\lambda_{\rm eff}=+0.10485824.}
\]
The point's SHA-256 is \texttt{476645438ea8a7ada42fea90e80174b772}\newline
\texttt{d45ce3e729d013d134e9848950c816}.
This is a local tree-level scalar benchmark, not global boundedness or a
viable GUT. The documented amendment changes a mixed coupling and compensates
its quadratic mass contribution, preserving the Higgs tuning while reducing
its heavy-singlet source. Direct-parent replay and a valley-polynomial
extraction check the positive quartic with shared tensor inputs. The failed
earlier point remains failed; this new witness excludes a universal claim
that all such Higgs tunings must have nonpositive quartics.

\newpage
\sect{4. Gauge matching: a measured mismatch, an incomplete verdict}
The recorded direct one-loop calculation matches the full parent to the SM
in a common $\overline{\rm MS}$ background-field Feynman-gauge convention,
with GUT-normalized $g_1$. It includes 290 heavy real scalar directions and
33 massive vector/Goldstone/ghost sets from the later positive-quartic point,
not the earlier failed tuning. The threshold ledger and matching program bind
that point's hash explicitly. Global boundedness remains unresolved.

Define $\alpha_i=g_i^2/(4\pi)$ and one-loop threshold functions $\lambda_i$ by
\[
\alpha_i^{-1}(\mu)=\alpha_{10}^{-1}(\mu)-\frac{\lambda_i(\mu)}{12\pi}.
\]
The complete state accounting gives
\[
\begin{gathered}
b_{\rm SM}=\left(\frac{41}{10},-\frac{19}{6},-7\right),\quad
b_S=\left(\frac{377}{30},\frac{77}{6},\frac{79}{6}\right),\quad
T_V=(8,6,5),\\
b_{{\rm SM},i}+b_{S,i}-\frac72 T_{V,i}=-\frac{34}{3}=b_{10},\qquad
\frac{d\lambda_i}{d\ln\mu}=6(b_{{\rm SM},i}-b_{10}).
\end{gathered}
\]
These are exact beta-jump and one-loop matching-scale checks. At
$\mu=\omega$, the frozen inverse-coupling inputs are
$(59.0100,29.5800,8.47458)$ at $M_Z$. The one-loop equations are
\[
\alpha_i^{-1}(M_Z)=\frac{4\pi}{g_{10}^2}
-\frac{\lambda_i^S+\lambda_i^V(g_{10})}{12\pi}
+\frac{b_i^{\rm SM}}{2\pi}\ln\frac{\omega}{M_Z}.
\]
They have no common perturbative solution for the frozen inputs. The recorded
least-squares residuals, in inverse-coupling units and channel order $(1,2,3)$,
are:
\begin{center}
\begin{tabular}{lrrrr}
\toprule Completed approximation & $r_1$ & $r_2$ & $r_3$ & $\max|r_i|$\\
\midrule One-loop running and thresholds & $0.36615$ & $-1.05980$ & $0.69389$ & $1.05980$\\
Two-loop SM running, one-loop thresholds & $0.26091$ & $-0.74548$ & $0.48479$ & $0.74548$\\
\bottomrule
\end{tabular}
\end{center}
The latter uses the SM running framework [5] and frozen reference boundaries
from [6], including top, bottom, and tau insertions with their one-loop
evolution. A separate inverse-coupling evolution with a different adaptive
integrator reproduces the residual to better than $2\times10^{-8}$ under
the frozen equations. Matching-scale variation $\mu/\omega\in[0.5,2]$
leaves maximum residuals from $0.74172$ to $0.74931$.

\sect{What the missing calculation must supply}
The exact two-loop RG logarithmic template has been derived, but its finite
model-dependent heavy constants have not. The published general matching
framework [4] provides methodology and a restricted comparator, not the finite
answer for these mixed, nondegenerate fields and multiple VEVs. The complete
Yukawa-dependent threshold also requires numerical family matrices and
right-handed-neutrino masses.

Small scale variation is not a rigorous bound on unknown finite constants.
Setting them to zero or freely tuning them would not compute a prediction.
Consequently the formal gauge verdict remains unresolved: neither unification
success nor exclusion has been established. The diagnostic fitted
$g_{10}\simeq0.5231$ and $\omega\simeq1.64\times10^{15}\,\mathrm{GeV}$
are not physical benchmark parameters. No replacement spectrum or flavor fit
is introduced in this note.

\newpage
\sect{5. Symmetry and quantum evidence: scope matters}
A separately named covariant physical gauge/scalar candidate satisfies exact
classical local gauge identities through factorized representation checks.
The recorded primary and tensor-functor replay compare the 328-real scalar
representation action and all 45 generator directions. The usual
$\delta(Dx)=-g_{10}\rho(\omega)Dx$ identity, invariant kinetic metric,
curvature identity, and invariant potential establish the classical statement.
This qualification is inactive; it does not replace a promoted component AST
or complete a quantum calculation.

Earlier literal component-action failures are retained. One preserved
two-field witness has nonzero gauge variation
\[
\delta\mathcal L_{\rm old}=
\frac{i g_{10}}{\sqrt2}(\partial_\mu C^4)
(Y\partial^\mu X-X\partial^\mu Y),
\]
with nonzero ghost Euler derivative. The later covariant candidate fixes the
kinetic structure by derivation. A failed encoded candidate is evidence about
that representation of the action, not a theorem against the gauge group.

Separately, the recorded one-loop parent scalar-pole projection has exact
rank $4+4+26=34$ and zero out-of-basis residual in its tested operator space.
Hermiticity and PQ selection-rule audits also pass within that scope. This
bounded radiative-closure result is not a complete anomaly or quantum
consistency theorem; it does not evaluate finite two-loop thresholds.

In the gauge-fixing/source investigation, a preserved normalization candidate
misses a selected auxiliary--antighost--ghost cubic pole. At one recorded
real-component witness, in common
$\hbar g_{10}\xi^2/(32\pi^2\epsilon)$ units, the coefficients are
\[
G_{\rm actual}=\frac{48\sqrt{10}}{25},\quad
G_{\rm candidate}=\frac{8\sqrt{10}}5,\quad
G_{\rm actual}-G_{\rm candidate}=\frac{8\sqrt{10}}{25}\ne0.
\]
Later bounded restoration investigations do not automatically erase this failed
candidate or establish full quantum consistency. The ledger remains at 34 of
39 executed counterterm items, with the relevant quantum gate blocked. New
uncommitted vector/ghost work was undergoing qualification at recovery and is
excluded from this note's earned claims. No completion is inferred from a
handoff or a live worker.

\sect{6. Remaining obligations and questions}
The later positive-quartic point establishes a bounded local tree-level scalar
witness. Physical viability still requires controlled interactions and loop
corrections; global boundedness and vacuum competition; a complete
consistent quantum gauge/source normalization, including measure, anomaly, and
infrared obligations; finite heavy-threshold matching; a fitted matter/neutrino
sector; and derived phenomenology. PQ anomaly coefficients, the physical axion,
global quotient, and topology also require their own analysis. The presence of
PQ selection rules alone does not establish an axion prediction or a strong-CP
solution for this variant. No experimental agreement is claimed.

Two useful discussion questions are: which analytic constraints on
$j^T\mathcal H^+j$ best screen other scalar points after the failed Higgs
tuning; and which controlled finite-threshold calculation or bound is needed
before assessing the frozen gauge mismatch? These questions do not presume
that missing corrections will rescue the model.

\sect{AI assistance and independence}
LLMs assisted calculation development, code, verification work, and preparation
of this note. Scientific direction and acceptance belong to the human author.
Distinct derivations and negative controls are useful, but shared action inputs,
some tensor/arithmetic code, Python/SymPy, and analyst limit independence.
Selected numerical work also uses NumPy, SciPy, mpmath, CVXPY/Clarabel; selected
tensor checks have a Julia replay. No external scientist's verification,
proof-assistant certification of the whole model, or endorsement is claimed.

\newpage
\sect{7. Evidence and reproduction map}
This note synthesizes recorded calculations. The reproduction guide separately
records package-byte checks and short read-only replays performed for publication.
Its input snapshot is local Git commit
\begin{center}\texttt{8a07e12f61350e1c6f488d0c52d7d23c676e156c}.\end{center}
The source scientific authority is \pathref{project_state.json}, version 113.
Public manifests authenticate exported bytes without promoting that ledger.
Historical custody remains unresolved, not a global pass.

Evidence paths below are relative to the source snapshot. The original Git
history is not included. The reproduction guide provides the packaged inputs,
scripts, dependencies, and evidence records for the note's claims.

{\small
\begin{tabular}{p{0.20\linewidth}p{0.74\linewidth}}
\toprule Claim & Evidence locator\\
\midrule
Scalar definition & \pathref{calculations/canonical_so10_scalar_reconstruction/}\newline
\pathref{PARENT_ACTION_V1.md}, \pathref{RESULT_V2.md}, \pathref{RECORD.md}\\[3pt]
Hessian and points & \pathref{calculations/canonical_so10_full_hessian/RESULT.md}\newline
\pathref{calculations/canonical_so10_scalar_benchmark/RESULT.md},\newline
\pathref{STAGE1_POINT.json}, \pathref{STAGE2_POINT.json}\\[3pt]
Failed / positive quartics & \pathref{calculations/canonical_so10_light_higgs_quartic/RESULT.md}\newline
\pathref{calculations/canonical_so10_positive_higgs/RESULT.md}, \pathref{POINT.json}\\[3pt]
Gauge residuals & \pathref{calculations/canonical_so10_direct_gauge_matching/RESULT.md}\newline
\pathref{calculations/canonical_so10_direct_gauge_next_order/RESULT.md},\newline
\pathref{next_order.json}\\[3pt]
Missing finite boundary & \pathref{calculations/canonical_so10_direct_two_loop_heavy_threshold/}\newline
\pathref{RESULT.md}, \pathref{YUKAWA_CONVENTION_V1.md}\\[3pt]
Scalar pole closure & \pathref{tools/two_loop_qft_compiler_v1/LAYER5B_M05_M08_RESULT.md}\\[3pt]
Classical and normalization & Under \pathref{tools/two_loop_qft_compiler_v1/}\newline
\pathref{m08_independent_replay_package_v1/}:\newline
\pathref{A10_PHYSICAL_ACTION_RECONCILIATION_RESULT.md},\newline
\pathref{A10_E5_AUXILIARY_VERTEX_RESULT.md}\\
\bottomrule
\end{tabular}
}

For targeted reproduction, the result files specify full prerequisites and
controls. From an isolated snapshot, representative entry points include
\begin{quote}\small\ttfamily
python -B calculations/canonical\_so10\_scalar\_reconstruction/verify\_action\_inventory.py\\
python -B calculations/canonical\_so10\_scalar\_benchmark/replay\_candidate\_exact.py\\
python -B calculations/canonical\_so10\_scalar\_benchmark/replay\_stage2\_exact.py\\
python -B calculations/canonical\_so10\_light\_higgs\_quartic/compute.py\\
python -B calculations/canonical\_so10\_positive\_higgs/verify\_exact\_point.py\\
python -B calculations/canonical\_so10\_direct\_gauge\_next\_order/independent\_replay.py
\end{quote}
Some calculation scripts write outputs. They should not run in a checkout with
an active immutable-input validation. The full Hessian, quartic, and quantum-pole
calculations were not rerun for publication, and no new parameter search was performed.

\sect{References}
{\small\setlength{\parskip}{3pt}
[1] K. S. Babu and S. Khan, \emph{A Minimal Non-Supersymmetric SO(10) Model:
Gauge Coupling Unification, Proton Decay and Fermion masses},
Phys. Rev. D \textbf{92}, 075018 (2015).
\href{https://arxiv.org/abs/1507.06712}{arXiv:1507.06712}.

[2] R. D. Peccei and H. R. Quinn, \emph{CP Conservation in the Presence of
Pseudoparticles}, Phys. Rev. Lett. \textbf{38}, 1440 (1977).
\href{https://doi.org/10.1103/PhysRevLett.38.1440}{doi:10.1103/PhysRevLett.38.1440}.

[3] A. Ernst, A. Ringwald and C. Tamarit, \emph{Axion Predictions in
SO(10) $\times$ U(1)$_{\PQ}$ Models}, JHEP \textbf{02} (2018) 103.
\newline\href{https://arxiv.org/abs/1801.04906}{arXiv:1801.04906}.

[4] W. Martens, \emph{Towards a Two-Loop Matching of Gauge Couplings in
Grand Unified Theories}, JHEP \textbf{01} (2011) 104.
\href{https://arxiv.org/abs/1011.2927}{arXiv:1011.2927}.

[5] M. Luo and Y. Xiao, \emph{Two-loop Renormalization Group Equations in
the Standard Model}, Phys. Rev. Lett. \textbf{90}, 011601 (2003).
\href{https://arxiv.org/abs/hep-ph/0207271}{arXiv:hep-ph/0207271}.

[6] D. Buttazzo et al., \emph{Investigating the near-criticality of the Higgs
boson}, JHEP \textbf{12} (2013) 089.
\href{https://arxiv.org/abs/1307.3536}{arXiv:1307.3536}.
}
\end{document}
